% Generated draft: requires mathematical and source review.
\documentclass[11pt]{article}
\usepackage[margin=25mm]{geometry}
\usepackage{fontspec}
\setmainfont{Latin Modern Roman}
\usepackage{amsmath,amssymb,mathtools,unicode-math}
\setmathfont{Latin Modern Math}
\usepackage{enumitem,needspace,xurl,hyperref}
\setlength{\parindent}{0pt}
\setlength{\parskip}{5pt}
\setlength{\emergencystretch}{3em}
\newcommand{\N}{\mathbb{N}}
\newcommand{\Z}{\mathbb{Z}}
\newcommand{\Q}{\mathbb{Q}}
\newcommand{\R}{\mathbb{R}}
\newcommand{\C}{\mathbb{C}}
\newcommand{\F}{\mathbb{F}}
\newcommand{\A}{\mathbb{A}}
\newcommand{\PP}{\mathbb{P}}
\newcommand{\E}{\mathbb{E}}
\newcommand{\eps}{\varepsilon}
\newcommand{\dd}{\,\mathrm{d}}
\begin{document}
\section*{UnsolvedMath ID: 30000078}
\textbf{Source problem code:} catalogue entry
\subsection{Definitions and mathematical statement}
Use the standard graded interpretation: let \(S=k[x_1,\ldots,x_m]\) be a polynomial ring with its fixed displayed coordinates, and let \(I\subsetneq S\) be a homogeneous ideal from the source's intended class of ``classical ideals.'' Calling \(I\) Gorenstein means that the quotient ring \(S/I\) is Gorenstein. For a monomial term order \(\prec\), set
\[
H=\operatorname{in}_{\prec}(I)
=(\operatorname{in}_{\prec}(f):f\in I).
\]
If \(J\) is a homogeneous ideal and \(\mathfrak m=(x_1,\ldots,x_m)\), its number of minimal homogeneous generators is
\[
\mu(J)=\dim_k(J/\mathfrak m J).
\]
The question is whether every Gorenstein classical ideal \(I\), in the intended family and field range, admits a term order \(\prec\) such that
\[
S/\operatorname{in}_{\prec}(I)\text{ is Gorenstein},
\qquad
\mu(\operatorname{in}_{\prec}(I))=\mu(I).
\]
The term order may depend on \(I\). The coordinates must remain fixed; allowing a preliminary linear coordinate change would be a different question. Equality is for minimal numbers of generators, not for the sizes of arbitrarily chosen generating lists, and does not by itself assert equality of all graded Betti numbers. Squarefreeness of the initial ideal is not required in the extracted statement.

The source does not define ``classical,'' specify the ground field or characteristic, or explicitly state homogeneity or whether initial ideals are taken with respect to term orders rather than weights. The display is therefore the conventional graded monomial-initial-ideal reading, with the universal family still to be specified from the report.

The supplied literature assessment checked 2026-08-21 reports positive constructions in Gorenstein Hibi, Segre, Veronese, determinantal, and Pfaffian families, but says no general theorem for every Gorenstein classical ideal in fixed coordinates was found. It does not reproduce the parameter restrictions or exact generator-preservation statements in those families; no all-parameter theorem for any listed family is inferred here.
\subsection{Short English statement}
In fixed polynomial coordinates, does every intended Gorenstein classical ideal admit a Gorenstein monomial initial ideal without increasing its minimal number of generators? The class of classical ideals and the exact ranges of known constructions need source verification.
\subsection{Sources}
\begin{itemize}
\item Local UnsolvedMath registry, permanent ID 30000078, catalogue entry, \emph{Gorenstein Initial Ideals with Minimal Generators}.
\item \emph{Combinatorial Commutative Algebra} (2005), Oberwolfach Reports, pp. 1703--1746, DOI \url{https://doi.org/10.4171/owr/2004/32}.
\item Conca--Hosten--Thomas, supplied label \emph{Nice Initial Complexes}: \url{https://arxiv.org/abs/math/0512283}.
\item Bruns--R\"omer, supplied label \emph{Gorenstein initial ideals}: \url{https://cel.hal.science/cel-00374634/document}.
\item Supplied literature triage checked 2026-08-21; supplied statement verification reports checking the official report. These source checks are not repeated here.
\end{itemize}
\end{document}
